% Part: first-order-logic
% Chapter: beyond
% Section: second-order-logic

\documentclass[../../../include/open-logic-section]{subfiles}

\begin{document}

\olfileid{fol}{byd}{sol}

\olsection{Tweede-ordelogica}

De taal van de tweede-ordelogica maakt het mogelijk niet alleen over
!!a{domain} van individuen te kwantificeren, maar ook over relaties op dat
!!{domain}. Gegeven een eerste-ordetaal $\Lang{L}$ voegen we voor elke $k$
!!{variable}s $R$ toe die over $k$-plaatsige relaties lopen, en staan we
kwantificatie over die variabelen toe. Als $R$ !!a{variable} voor een
$k$-plaatsige relatie is en $t_1$, \dots,~$t_k$ gewone (eerste-orde) termen
zijn, dan is $\Atom{R}{t_1,\dots,t_k}$ een atomaire !!{formula}. Verder wordt
de verzameling !!{formula}s net als bij de eerste-ordelogica gedefinieerd,
met aanvullende clausules voor tweede-ordekwantificatie. Merk op dat we alleen
!!{identity} voor eerste-ordetermen hebben: als $R$ en $S$ relatie!!{variable}s
van dezelfde plaatsigheid~$k$ zijn, kunnen we $\eq[R][S]$ definiëren als
afkorting voor
\[
\lforall[x_1 \dots][\lforall[x_k][(\Atom{R}{x_1, \dots, x_k} \liff
  \Atom{S}{x_1, \dots, x_k})]].
\]

De regels voor tweede-ordelogica breiden de kwantorregels eenvoudig uit tot
de nieuwe tweede-ordevariabelen. Daarbij moet echter zorgvuldig worden
uitgelegd hoe deze variabelen samenhangen met de !!{predicate}s van
~$\Lang{L}$ en, algemener, met de !!{formula}s van~$\Lang{L}$. Relatievariabelen
gelden in elk geval als termen, zodat we gevolgtrekkingen hebben van de vorm
\[
!A(R) \Proves \lexists[R][!A(R)]
\]
Maar als $\Lang{L}$ de taal van de rekenkunde is, met het constante
relatiesymbool~$<$, verwachten we ook dat de volgende gevolgtrekking geldig is:
\[
x < y \Proves \lexists[R][\Atom{R}{x,y}]
\]
of, voor een gegeven !!{formula}~$!A$,
\[
!A(x_1, \dots, x_k) \Proves \lexists[R][\Atom{R}{x_1,\dots,x_k}]
\]
Meer in het algemeen zouden we gevolgtrekkingen van de volgende vorm kunnen
toestaan:
\[
\Subst{!A}{\lambd[\vec x][!B(\vec x)]}{R} \Proves
\lexists[R][!A]
\]
Hier staat $\Subst{!A}{\lambd[\vec x][!B(\vec x)]}{R}$ voor het resultaat
van de vervanging van elke atomaire !!{formula} van de vorm
$\Obj{R}{t_1,\dots,t_k}$ in~$!A$ door $!B(t_1, \dots, t_k)$. Deze laatste
regel komt overeen met een {\em comprehensieschema}, dat wil zeggen een
axioma van de vorm
\[
\lexists[R][\lforall[x_1, \dots, x_k][(!A(x_1, \dots, x_k) \liff
\Atom{R}{x_1, \dots, x_k})]],
\]
één voor elke !!{formula}~$!A$ in de tweede-ordetaal waarin $R$ niet vrij
voorkomt. (Opgave: toon aan dat dit schema inconsistent is als $R$ wel in
~$!A$ mag voorkomen!)

Wanneer logici spreken over de ``axioma's van de tweede-ordelogica'', bedoelen
ze gewoonlijk de minimale uitbreiding van de eerste-ordelogica met
tweede-ordekwantorregels en het comprehensieschema. Het is echter vaak
interessant zwakkere subsystemen van deze axioma's en regels te bestuderen.
Het axiomaschema voor comprehensie is in zijn volle algemeenheid bijvoorbeeld
\emph{impredicatief}: het maakt het mogelijk het bestaan te poneren van een
relatie $\Atom{R}{x_1, \dots, x_k}$ die wordt ``gedefinieerd'' door
!!a{formula} met tweede-ordekwantoren. Die kwantoren lopen over de verzameling
van al zulke relaties, en die verzameling bevat $R$ zelf!{} Rond de overgang
naar de twintigste eeuw was een veel voorkomende reactie op de paradox van
Russell om zulke definities verantwoordelijk te stellen en ze bij de
grondslagen van de wiskunde te vermijden. Als tweede-ordekwantoren in de
!!{formula}~$!A$ worden verboden, krijgen we een {\em predicatieve} vorm van
comprehensie, die iets zwakker is.

Semantisch kunnen we een tweede-orde!!{structure} opvatten als een
eerste-orde!!{structure} voor de taal, samen met een verzameling relaties op
het !!{domain} waarover de tweede-ordekwantoren lopen (preciezer: voor elke
$k$ is er een verzameling relaties met plaatsigheid~$k$). Als comprehensie in
het !!{derivation} systeem is opgenomen, geldt bovendien de eis dat het
``tweede-ordedeel'' genoeg relaties bevat om aan de comprehensieaxioma's te
voldoen---anders is het !!{derivation} systeem niet correct!{} Een eenvoudige
manier om genoeg relaties te garanderen is het tweede-ordedeel uit
\emph{alle} relaties op het eerste-ordedeel te laten bestaan. Zo'n
!!a{structure} heet \emph{volledig} en is in zekere zin werkelijk de
``bedoelde !!{structure}'' voor de taal. Als we ons tot volledige
!!{structure}s beperken, krijgen we de zogeheten \emph{volledige}
tweede-ordesemantiek. Om zo'n !!{structure} vast te leggen volstaat het dan
het eerste-ordedeel vast te leggen, want daaruit volgt impliciet wat het
tweede-ordedeel bevat.

Samengevat is het enigszins dubbelzinnig wat met tweede-ordelogica wordt
bedoeld. Bij het !!{derivation} systeem kan het gaan om
\begin{enumerate}
\item een ``minimaal'' tweede-orde!!{derivation} systeem, met enkele
  comprehensieaxioma's;
\item het ``standaard'' tweede-orde!!{derivation} systeem, met volledige
  comprehensie.
\end{enumerate}
Bij de semantiek kan het gaan om
\begin{enumerate}
\item de ``zwakke'' semantiek, waarin !!a{structure} uit een eerste-ordedeel
  en een tweede-ordedeel bestaat dat groot genoeg is om aan de
  comprehensieaxioma's te voldoen;
\item de ``standaardsemantiek'' voor tweede-ordelogica, waarin alleen
  volledige !!{structure}s worden beschouwd.
\end{enumerate}
Als logici niet vermelden welk !!{derivation} systeem of welke semantiek ze
bedoelen, doelen ze gewoonlijk op het tweede geval in beide lijsten. Het
voordeel van deze semantiek is dat ze, zoals we zullen zien, categorische
beschrijvingen van veel natuurlijke wiskundige structuren oplevert. Het
!!{derivation} systeem is bovendien behoorlijk sterk en correct voor deze
semantiek. Het nadeel is dat het !!{derivation} systeem \emph{niet} volledig
is voor de semantiek; sterker nog, \emph{geen} effectief gegeven
!!{derivation} systeem is volledig voor de volledige tweede-ordesemantiek.
Voor de verzwakte semantiek is het !!{derivation} systeem daarentegen
\emph{wel} volledig. Dit houdt in dat er, als een zin niet bewijsbaar is,
\emph{een of andere} !!{structure} bestaat waarin die zin onwaar is, al hoeft
dat niet de volledige structuur te zijn.

De taal van de tweede-ordelogica is zeer expressief. Eenplaatsige relaties
kunnen met deelverzamelingen van het !!{domain} worden geïdentificeerd, zodat
we in het bijzonder over die verzamelingen kunnen kwantificeren. Zo kan
inductie voor de natuurlijke getallen met één enkel axioma worden uitgedrukt:
\[
\lforall[R][((\Atom{R}{\Obj{0}} \land \lforall[x][(\Atom{R}{x} \lif
    \Atom{R}{x'})]) \lif \lforall[x][\Atom{R}{x}])].
\]
Als de taal van de rekenkunde de symbolen $\Obj 0, \Obj \prime, +,
\times$ en $<$ bevat, kunnen we de volgende axioma's toevoegen om hun gedrag
te beschrijven:
\begin{enumerate}
\item $\lforall[x][\lnot x' = \Obj 0]$
\item $\lforall[x][\lforall[y][(s(x) = s(y) \lif x = y)]]$
\item $\lforall[x][(x + \Obj 0) = x]$
\item $\lforall[x][\lforall[y][(x + y') = (x + y)']]$
\item $\lforall[x][(x \times \Obj 0) = \Obj 0]$
\item $\lforall[x][\lforall[y][(x \times y') = ((x \times y) + x)]]$
\item $\lforall[x][\lforall[y][(x < y \liff \lexists[z][y = (x + z')])]]$
\end{enumerate}
Het is niet moeilijk aan te tonen dat deze axioma's samen met het bovenstaande
inductieaxioma een categorische beschrijving geven van de
!!{structure}~$\Struct{N}$, het standaardmodel van de rekenkunde, mits we de
volledige tweede-ordesemantiek gebruiken. Neem een willekeurige
!!{structure}~$\Struct{M}$ waarin deze axioma's waar zijn. Definieer een
functie~$f$ van $\Nat$ naar het !!{domain} van $\Struct{M}$ met gewone
recursie op~$\Nat$, zodat $f(0) = \Assign{\Obj 0}{M}$ en
$f(x+1) = \Assign{\prime}{M}(f(x))$. Uit gewone inductie op~$\Nat$ en het
feit dat axioma's (1) en~(2) in~$\Struct M$ gelden, volgt dat $f$
!!{injective} is. Om te zien dat $f$ !!{surjective} is, nemen we voor~$P$ de
verzameling elementen van~$\Domain{M}$ die in het bereik van~$f$ liggen.
Omdat $\Struct M$ volledig is, ligt $P$ in het tweede-orde!!{domain}. Uit de
constructie van~$f$ weten we dat $\Assign{\Obj 0}{M}$ in~$P$ ligt en dat
$P$~gesloten is onder $\Assign{\prime}{M}$. Omdat het inductieaxioma in
$\Struct M$ geldt (in het bijzonder voor~$P$), is $P$~gelijk aan het hele
eerste-orde!!{domain} van~$\Struct M$. Daarmee is $f$ !!a{bijection}. Met
Even eenvoudig is met herhaalde gewone inductie aan te tonen dat $f$ een
homomorfisme is op $\Nat$.

Verzamelingstheoretisch is een functie slechts een bijzonder soort relatie.
Een eenplaatsige functie~$f$ kan bijvoorbeeld worden geïdentificeerd met een
tweeplaatsige relatie~$R$ die voldoet aan
$\lforall[x][\lexists![y][R(x,y)]]$. We kunnen dus ook over functies
kwantificeren. Met de volledige semantiek kunnen we vervolgens de klasse van
oneindige !!{structure}s definiëren als de klasse van !!{structure}s
$\Struct M$ waarvoor een injectieve functie bestaat van het !!{domain} van
$\Struct M$ naar een echte deelverzameling van zichzelf:
\[
\lexists[f][(\lforall[x][\lforall[y][(\eq[f(x)][f(y)] \lif \eq[x][y])]]
  \land \lexists[y][\lforall[x][\eq/[f(x)][y]]])].
\]
De negatie van deze zin definieert dan de klasse van eindige !!{structure}s.

Ook kunnen we de klasse van welordeningen definiëren door aan de definitie
van een lineaire ordening het volgende toe te voegen:
\[
\lforall[P][(\lexists[x][\Atom{P}{x}] \lif \lexists[x][(\Atom{P}{x}
    \land \lforall[y][(y < x \lif \lnot \Atom{P}{y})])])].
\]
Dit zegt dat elke niet-lege verzameling een kleinste element heeft, waarbij
we een ``verzameling'' met een eenplaatsige relatie identificeren. Als ander
voorbeeld kunnen we de samenhangendheid van grafen uitdrukken door te zeggen
dat de punten niet op een niet-triviale manier in onderling niet-verbonden
delen kunnen worden gesplitst:
\[
\lnot \lexists[A][(\lexists[x][A(x)] \land \lexists[y][\lnot A(y)]
  \land \lforall[w][\lforall[z][((\Atom{A}{w} \land \lnot \Atom{A}{z})
      \lif \lnot \Atom{R}{w,z})]])].
\]
Als verdere opgave kan men proberen de klasse te definiëren van eindige
!!{structure}s waarvan het !!{domain} een even aantal elementen heeft. Nog
opmerkelijker is dat de reële getallen categorisch kunnen worden beschreven
als een volledig geordend lichaam dat de rationale getallen bevat.

Kortom, tweede-ordelogica is veel expressiever dan eerste-ordelogica. Dat is
het goede nieuws; nu het slechte. We hebben al vermeld dat er geen effectief
!!{derivation} systeem bestaat dat volledig is voor de volledige
tweede-ordesemantiek. Veel eigenschappen van de eerste-ordelogica ontbreken,
ten goede of ten kwade, waaronder compactheid en de stellingen van
L\"owenheim--Skolem.

Als we daarentegen bereid zijn de volledige tweede-ordesemantiek voor de
zwakkere op te geven, is het minimale tweede-orde!!{derivation} systeem
volledig voor deze semantiek. Als we $\Proves$ lezen als ``bewijsbaar in het
minimale systeem'' en $\Entails$ als ``volgt logisch in de zwakkere
semantiek'', kunnen we dus aantonen dat uit $\Gamma \Entails !A$ steeds
$\Gamma \Proves !A$ volgt. Willen we bepaalde comprehensieaxioma's in het
!!{derivation} systeem opnemen, dan moeten we de semantiek beperken tot
tweede-ordestructuren die aan deze axioma's voldoen. Als $\Delta$
bijvoorbeeld een verzameling comprehensieaxioma's is (eventueel alle), dan
geldt: als $\Gamma \cup \Delta \Entails !A$, dan
$\Gamma \cup \Delta \Proves !A$. In het bijzonder bestaat er, als $!A$ niet
bewijsbaar is met de comprehensieaxioma's die we beschouwen, een model van
$\lnot !A$ waarin die comprehensieaxioma's niettemin gelden.

De eenvoudigste manier om in te zien dat de volledigheidsstelling voor de
zwakkere semantiek geldt, is tweede-ordelogica als volgt als een meersoortige
logica op te vatten. Eén soort wordt geïnterpreteerd als het gewone
``eerste-orde'' !!{domain}; daarnaast hebben we voor elke $k$ !!a{domain} van
``relaties met plaatsigheid $k$.'' De taal bevat ingebouwde relatiesymbolen
``$\Atom{\Obj{true}_k}{R,x_1,\dots,x_k}$''. Daarmee wordt uitgedrukt dat $R$
geldt voor $x_1$, \dots,~$x_k$, waarbij $R$ een variabele van de soort
``$k$-plaatsige relatie'' is en $x_1$, \dots,~$x_k$ objecten van de
eerste-ordesoort zijn.

Met deze identificatie is de zwakke tweede-ordesemantiek in wezen de gewone
semantiek voor meersoortige logica. We hebben al gezien dat meersoortige
logica in eerste-ordelogica kan worden ingebed. Afgezien van de vertalingen
heen en terug is de zwakkere opvatting van tweede-ordelogica dus eigenlijk
eerste-ordelogica in vermomming: het !!{domain} bevat zowel ``objecten'' als
``relaties'', die aan de bijbehorende axioma's zijn onderworpen.

\end{document}
